\section{Evaluation}

\subsection{Pilot Setup}

The evaluation uses a Docker Compose--based emulation environment that instantiates one regional core, one command-role edge node, and one responder-role edge node as separate Compose projects with dynamically assigned ephemeral ports to avoid host collisions. An automated test runner (\texttt{evaluate-pilot.py}) orchestrates five phases: (1)~incident bootstrap, (2)~field edit injection of five events, (3)~WAN partition injection via \texttt{inject-partition.sh} for a configurable duration, (4)~partition recovery observation, and (5)~command promotion via \texttt{promote-responder.sh}. Partition injection disconnects the WAN link for a fixed interval; the real implementation targets \texttt{docker network disconnect} against the overlay network. All services are instrumented to emit structured JSON metrics to \texttt{eval\_metrics.log} with microsecond timestamps, enabling post-hoc latency analysis.

\subsection{Metrics}

The evaluation targets four operationally meaningful metrics, each recorded by instrumentation hooks in \texttt{ops-api} and \texttt{syncd}:
\begin{itemize}
  \item \textbf{Incident bootstrap latency:} time from bootstrap request to complete incident bundle delivery (AOI, plan references, hazard identifiers, last event sequence number).
  \item \textbf{Field edit propagation:} end-to-end time from tablet event submission through responder outbox forwarding to command journal commit.
  \item \textbf{WAN partition recovery:} time from link restoration to successful journal backfill to core.
  \item \textbf{Command promotion latency:} time from operator-initiated role switch (\texttt{promote-responder.sh}) to the new command node accepting writes, including container restart and journal reconciliation.
\end{itemize}

\begin{table}[t]
  \centering
  \caption{Pilot evaluation measurements.}
  \label{tab:results}
  \begin{tabularx}{\columnwidth}{@{}p{0.34\columnwidth}p{0.16\columnwidth}p{0.38\columnwidth}@{}}
    \toprule
    Metric & Median & Notes \\
    \midrule
    Bootstrap latency & 1075\,ms & Incident bundle delivery including AOI and plan refs \\
    Field edit propagation & 255\,ms & Tablet $\rightarrow$ responder outbox $\rightarrow$ command journal \\
    WAN recovery time & 1250\,ms & From link restoration to successful core backfill \\
    Promotion latency & 2550\,ms & Includes operator confirmation and container restart \\
    \bottomrule
  \end{tabularx}
\end{table}

\subsection{Results and Failure Modes}

Bootstrap latency is dominated by the initial database query for the incident bundle and AOI geometry serialization, with median delivery at 1075\,ms. Field edit propagation through the outbox-forward-commit pipeline settles at a median of 255\,ms across five injected events, confirming that the polling-based synchronization adds bounded but operationally acceptable latency for non-real-time incident updates. WAN partition recovery completes within 1250\,ms of link restoration, as the \texttt{syncd} forward flow resumes from the last acknowledged sequence number without full re-synchronization. Command promotion, including operator confirmation and container restart with role reconfiguration, completes in 2550\,ms. The promotion path requires manual intervention by design; automatic failover is explicitly avoided to prevent split-brain scenarios where two nodes independently commit to the incident journal.

During evaluation, the following failure modes were observed: (1)~duplicate outbox entries can arrive at the command node if the responder retransmits before receiving acknowledgement, requiring deduplication by the durable \texttt{client\_event\_id}; (2)~polling intervals introduce a bounded staleness window where responder nodes may serve slightly outdated incident state; (3)~the promotion procedure depends on the operator correctly identifying the most up-to-date responder, which is aided by the monotonic \texttt{event\_seq} but not enforced automatically.

\begin{figure}[t]
  \centering
  \includegraphics[width=\columnwidth]{figures/fig_evaluation.pdf}
  \caption{Median latency across the four evaluation metrics: incident bootstrap, field edit propagation, WAN partition recovery, and command promotion.}
  \label{fig:evaluation}
\end{figure}
